\title{Nominal Budgets and Realized Resources in Closed-API Large-Language-Model Inference: A Performance Evaluation Protocol}

\author[aff1]{Soroush Vahidi\corref{cor1}\fnref{fn1}}
\ead{sv96@njit.edu}

\cortext[cor1]{Corresponding author}
\fntext[fn1]{Sole author. ORCID: \href{https://orcid.org/0000-0003-1934-6282}{0000-0003-1934-6282}.}

\address[aff1]{Ying Wu College of Computing, New Jersey Institute of Technology, Newark, New Jersey, USA}

\begin{abstract}
Closed-API large language model (LLM) evaluations often match systems by nominal inference budget,
but equal nominal budgets need not imply equal realized resources or fair attribution between answer
discovery and final-answer selection. We present a unified protocol for closed-API, budgeted
inference that records componentwise resources, evaluates identical-pool selector controls, labels
provider transfer for failure-mined rules, preserves protocol-blocked outcomes, and discloses
incomplete telemetry. We exercise it on 15 completed provider-by-dataset cells with $3{,}394$ paired
examples; Fireworks$\times$GPQA-Diamond is one protocol-blocked cell. On the shared four-generator
pool, identical-pool plurality (Pooled-4) exceeds the failure-trace gated selector (FTA), $66.53\%$
versus $65.00\%$ (McNemar $p{=}0.00027$). FTA's dominant gate transfers poorly in Azure mathematics
cells ($2/18$ and $7/31$ rescue/regression patterns). Resource reconstruction changes efficiency
interpretation: Frontier successful completions rise from $2.78$ to $5.38$ of nominal budget
$B{=}6$, while token and dollar fields remain lower bounds. Nominal budget alone is therefore
insufficient for defensible closed-API inference-performance conclusions.
\end{abstract}

\begin{keyword}
budgeted inference \sep performance measurement \sep resource accounting \sep closed-API inference \sep large language models
\end{keyword}
